% -----------------------------------------------
% スライド1: 表紙＋目標＋用語・定義
% -----------------------------------------------
\begin{frame}{\CourseName\quad 第\LectureNum 回「\LectureTitle 」}
  {\small \TermName\quad \TeacherName\quad ／\quad 教科書 \TextPages}

  \begin{block}{用語・定義}
  \begin{description}[labelwidth=5.5em]
    \item[\kw{関数}] $x$ の値が決まると $y$ の値がただ1つ定まる対応関係
    \item[\kw{独立変数}] 自由に値を取れる変数（入力側）$x$
    \item[\kw{従属変数}] 独立変数に従って値が決まる変数（出力側）$y$
    \item[\kw{定義域}] 独立変数 $x$ が取り得る値の集合（範囲）
    \item[\kw{値域}] 従属変数 $y$ が取り得る値の集合（範囲）
    \item[\kw{変域}] 変数のとる値の範囲
    \item[\kw{f(a)}] 関数 $f$ に $x=a$ を代入した値：$y=f(a)$
  \end{description}
  \end{block}
\end{frame}

% -----------------------------------------------
% スライド2: 座標平面
% -----------------------------------------------
\begin{frame}{座標平面}
  \begin{block}{用語・定義}
  \begin{description}[labelwidth=5.5em]
    \item[\kw{座標平面}] 2本の数直線（$x$ 軸・$y$ 軸）を直交させた平面
    \item[\kw{原点}] 2軸の交点 O（$x=0, y=0$）
    \item[\kw{座標}] 点の位置を表す組 $(x, y)$
  \end{description}
  \end{block}
\end{frame}

% -----------------------------------------------
% スライド3: 公式・性質
% -----------------------------------------------
\begin{frame}{公式・性質}
  \begin{block}{}
  \begin{itemize}
    \item 1次関数 $y=ax+b$
    \item 定数関数 $y=c$
    \item $x=d$（垂直な直線）
  \end{itemize}
  \end{block}
\end{frame}

% -----------------------------------------------
% スライド4: 例1→問1
% -----------------------------------------------
\begin{frame}{例1→問1}
  \begin{exampleblock}{例1}
    （板書で解説）
  \end{exampleblock}

  \vfill

  \begin{block}{問1\quad 自分で解いてみよう}
    （ノートに解くこと）
  \end{block}
\end{frame}

% -----------------------------------------------
% スライド5: 例2・例3
% -----------------------------------------------
\begin{frame}{例2・例3}
  \begin{exampleblock}{例2}
    （板書で解説）
  \end{exampleblock}

  \vfill

  \begin{exampleblock}{例3}
    （板書で解説）
  \end{exampleblock}
\end{frame}

% -----------------------------------------------
% スライド6: 問2
% -----------------------------------------------
\begin{frame}{問2}
  \begin{block}{問2\quad 自分で解いてみよう}
    （ノートに解くこと）
  \end{block}
\end{frame}

% -----------------------------------------------
% まとめ
% -----------------------------------------------
\begin{frame}{まとめ}
  \begin{itemize}
    \item 関数とは $x$ の値に $y$ がただ1つ対応する関係
    \item 定義域は $x$ の範囲、値域は $y$ の範囲
    \item 1次関数 $y=ax+b$ のグラフは傾き $a$・切片 $b$ の直線
  \end{itemize}

  \vfill
  {\small 問題集Basic：143, 144}
\end{frame}
